
This work investigated whether per-sample loss trajectories during training can identify samples affected by spurious correlations. On the Waterbirds benchmark, trajectory features predicted minority group membership with AUROC = 0.9452, and this signal showed differential attenuation under GroupDRO versus random reweighting, providing evidence for spurious-specificity. Initial loss alone achieved comparable performance (AUROC = 0.9473), suggesting detection is possible from the first training epoch. A secondary hypothesis about delayed curvature stabilization was not supported.

These results demonstrate that loss trajectory analysis can serve as a diagnostic for spurious correlation vulnerability on this benchmark. The approach requires only standard training with per-sample loss logging and could be integrated into existing training pipelines. Whether these findings generalize to other spurious correlation settings, other model architectures, or translate to successful intervention remain open questions for future work.
